problem:  
Let \((M^4,g)\) be a closed oriented Riemannian 4‑manifold with positive Yamabe invariant.  
Denote by \(Q_g\) the Branson \(Q\)‑curvature and by \(P_g\) the Paneitz operator, satisfying for each conformal metric \(\tilde g=e^{2u}g\):
\[
P_g u + 2Q_g = 2Q_{\tilde g} e^{4u}.
\]  
A conformal metric \(\tilde g\) is said to have **constant Q‑curvature** \(\bar Q\) if
\[
P_g u + 2Q_g = 2\bar Q e^{4u}.
\]

Assume \(P_g\) is **nonnegative with trivial kernel**, but not necessarily strictly positive.

**Problem:**  
Investigate the analytic–geometric threshold governing compactness and blow‑up of solutions \(u_i\) to the constant Q‑curvature equation in this borderline case. In particular,

1. Is there a conformally invariant **spectral–topological condition**—formulated via the first nonzero eigenvalue of \(P_g\), the sign distribution of its Green’s function, or a global curvature functional—that guarantees precompactness of the solution set \(\{u : P_g u + 2 Q_g = 2 \bar Q e^{4u}\}\) in \(C^\infty(M)\)?  
2. If compactness fails, describe the blow‑up mechanisms: do the measures \(Q_{g_i} dV_{g_i}\) (where \(g_i=e^{2u_i}g\)) converge weakly to sums of Dirac masses  
\[
Q_{g_i} dV_{g_i}\rightharpoonup \sum_j 8\pi^2 m_j \delta_{p_j}, \qquad m_j\in\mathbb N,
\]
and does this **quantization** of \(Q\)-curvature correspond to concentration of the Green’s function maxima of \(P_g\)?  
3. Does bubbling occur only when \(\int_M Q_{g_i}\,dV_{g_i} \to 8\pi^2\chi(M)\), equivalently when the Weyl energy \(\int_M |W_{g_i}|^2 dV_{g_i}\to 0\), meaning that blow‑up is intimately linked with **approaching conformal flatness**?

Why is it a "good" problem:  
This problem rigorously targets the most delicate and least understood part of the constant Q‑curvature program: the **borderline spectral regime** where the Paneitz operator loses strict positivity but retains a trivial kernel. The equation is structurally critical—conformally invariant of fourth order—and its compactness and blow‑up analysis would deepen understanding of geometric quantization phenomena in four dimensions.  

It is a “good” problem because it:  
- Addresses a **precise spectral threshold** for analysis and geometry of 4‑manifolds, which remains open;  
- Extends the known constant Q‑curvature compactness theory beyond the positive Paneitz case;  
- Links **spectral geometry** (sign and structure of the Green’s function of \(P_g\)), **nonlinear analysis** (compactness, bubbling, quantization), and **topology** (\(\chi(M)\), conformal flatness limits);  
- Mirrors the elegance of the 2‑D \(8\pi\) quantization phenomena in a higher‑dimensional, conformally invariant setting;  
- Retains simplicity in formulation—compactness versus blow‑up for one geometric PDE—while probing the deep analytic and geometric architecture of Q‑curvature on 4‑manifolds.